\section{Model Extraction Attacks}
Machine Learning as a Service (MLaaS) creates a conflict between providing high-precision results through rich APIs and maintaining model confidentiality for monetization. Models are valuable because training is resource-expensive, involves novel techniques, or uses sensitive/exclusive data.

\definition{Model Extraction Attack}{An attack (usually grey-box) against MLaaS where an adversary aims to "steal" a remote ML model. The attacker queries the target model \f{F_{\theta}} with a dataset \f{Z} to learn a local surrogate model \f{g_{\hat{\theta}}}. The main challenge is minimizing the number of queries required.}

\subsection{Attack Objectives}
\begin{enumerate}
\item \b{High-level objectives}: Economic incentives (avoiding query charges, copying a business model) or facilitating further attacks (adversarial examples, inference attacks).
\item \b{Task-accuracy extraction}: Maximizing the expected task accuracy 
\cf{\max(E_{(x,y)\sim P_{xy}}[g_{\hat{\theta}}(x)=y])}
\item \b{Fidelity extraction}: Maximizing the behavioral similarity to the target model 
\cf{\max(E_{x\sim P_{x}}[sim(g_{\hat{\theta}}(x),F_{\theta}(x))])}
\item \b{Functionally-equivalent extraction}: Achieving exact identical outputs \f{g_{\hat{\theta}}(x)=F_{\theta}(x)\ \forall x \in X}
\item \b{Exact copy}: Reconstructing the exact prediction function \f{g=F} and parameters \f{\hat{\theta}=\theta}.
\item Other objectives: Extracting hyper-parameters, the model architecture, or specific learned parameters.
\end{enumerate}



\subsection{Attacker Capabilities}
MEAs are rarely pure black-box attacks. More often than not, the attacker has at least some form of knowledge about the target, as well as different types of data availability.
\begin{enumerate}
\item[D1.] \b{Abundant Data}: Access to a large, representative dataset.
\item[D2.] \b{Partial Data}: Access to only a small representative dataset.
\item[D3.] \b{No Data}: No target-domain data available, forcing the attacker to synthesize query data.
\item[]
\item[K1.] \b{System Knowledge}: Knowledge of the learning algorithm, hyper-parameters, and the optimizer.
\item[K2.] \b{Task-domain Knowledge}: Knowledge of the task \f{T} and implicitly the data distribution \f{D}.
\item[K3.] \b{Output Knowledge}: Defined by the API interface (e.g., returning hard labels, softmax scores, or top-k probabilities).
\item[K4.] \b{Model-view Knowledge}: Additional information such as assumed architecture, gradient information from XAI outputs, or side-channel information.
\end{enumerate}



\subsection{Attack Strategies}
\begin{itemize}
\item \b{Analytic Attacks}: Deriving the surrogate model through an analytic procedure (e.g., solving equations). Requires detailed information about the target (e.g., deriving bias and weights for a Linear SVM by carefully crafting inputs).\\[.5em]
\b{Example}: For a linear SVM (\f{f(z)=zw^T+b}) an attacker can derive the bias ($b$) by querying $z=(0,...0)$ and then repeat for each dimension to find the weights ($w$) with \f{z=(0,1,0,...0)}.

\item \b{Learning-based Attacks}: Querying input-output pairs to actively train a surrogate model.
\end{itemize}



\subsection{Connections to other ML Concepts}
MEA is highly connected to other ML fields, often reusing their techniques to refine the surrogate model.

\subsubsection{Active Learning (AL)}
AL concepts are used to minimize labeling effort by querying the target model only for the most informative or uncertain samples (e.g. ActiveThief).

\subsubsection{White-Box AdvEx (PRADA)}
Used to generate synthetic adversarial pairs that accurately describe the decision boundary of the surrogate (e.g. PRADA using FGSM with partial data).

\subsubsection{Black-Box AdvEx (MAZE)}
A data-free model stealing attack that uses a generator to create synthetic data. A generator loss \f{L_G} maximizes the disagreement between the target and the surrogate, while a surrogate loss \f{L_S} maximizes their agreement. Because calculating \f{L_G} requires gradients from the target model, MAZE uses \b{Zeroth-Order Optimization} (gradient estimation).

\subsubsection{Knowledge-Distillation (DisGUIDE)}
Training a small student model (surrogate) based on the outputs of a large teacher model (target).\\[.5em]
\b{Zero-shot KD}: Relates to "No-data ME" where training samples are synthesized with a generator.\\[.5em]
\b{DisGUIDE}: A no-data extraction attack that improves query efficiency by using multiple student surrogates. It generates synthetic query samples specifically targeting regions where the two student surrogates disagree.



\subsection{Defenses}
Defending against model extraction is notoriously difficult because it is hard to quantify how much the adversary has already learned.

\begin{itemize}
\item \b{Stateful Application}: Monitoring access to the prediction function to detect unusual querying behavior. This is highly vulnerable to Sybil attacks.
\item \b{Watermarking}: Ownership verification by introducing unique, robust patterns into the model.
\begin{itemize}
    \item \b{Black-box verification:} Backdooring the model so that a stolen surrogate will also respond to a specific set of secret adversarial examples (probabilistic result only).
    \item \b{White-box verification:} Encoding patterns directly into the parameters.
\end{itemize}
\end{itemize}





